\section*{Abstract}

This project presents an Artificial Intelligence system capable of selecting discard actions in Riichi Mahjong, a traditional Japanese game of imperfect information.

Using a Transformer-based architecture, the model is trained through imitation learning on a dataset composed of real games collected from Tenhou.net. The model is evaluated in terms of predictive accuracy, and the implications of the obtained results are discussed.

In addition to conventional accuracy metrics, this work proposes a complementary evaluation method referred to as \textit{Strategic Accuracy}. It uses categories defined through Shanten and Ukeire to analyse hand efficiency beyond exact agreement with the observed action.

On a reserved evaluation set containing 52,563 game states, the model trained on 500,000 states achieves a top-1 accuracy of 65.64\%, a top-3 accuracy of 92.47\%, and a strategic accuracy of 94.01\%. Compared with the checkpoint identified as 100k, the second model improves all predictive metrics and reduces Shanten-degrading decisions from 5.03\% to 1.59\%. These results indicate that the model learns relevant discard-selection patterns, although the proposed strategic evaluation remains limited to the aspects represented by Shanten and Ukeire.

\vspace{0.5cm}

\textbf{Keywords:}

Riichi Mahjong,
Artificial Intelligence,
Transformers,
Imitation Learning,
Strategic Accuracy,
Model Evaluation,
Imperfect Information Games.
